\section{可复现算法与逐日流水}

\subsection{核心算法伪代码}

\begin{resultbox}{算法A：确定天气时空网络}
\begin{enumerate}[leftmargin=1.8em]
  \item 从场景对象读取邻接图、天气、特殊节点与资源参数；
  \item 创建日期—位置二元流变量、行动变量、购买量、库存和现金变量；
  \item 加入每日唯一行动、位置流、沙暴禁行、矿山、村庄、资源、现金和负重约束；
  \item 最大化终点现金与半价回收，调用HiGHS；
  \item 把最优变量转换为逐日计划，交独立模拟器回放；任何规则不一致则拒绝输出。
\end{enumerate}
\end{resultbox}

\begin{resultbox}{算法B：天气先揭示的有限时域MDP}
\begin{enumerate}[leftmargin=1.8em]
  \item 枚举所有满足现金与负重的整数初始采购；
  \item 从截止日反向计算 $V_t(t,l,w,f)$，对当天每种支撑天气先观察再选行动；
  \item 若任一正概率后继无法按时到达，则当前动作不可行；
  \item 在第0日选择期望值最大的采购，并保存确定性策略表；
  \item 对全部天气支撑逐条回放，复核Bellman值和失败数。
\end{enumerate}
\end{resultbox}

\begin{resultbox}{算法C：第五关完整最佳响应}
\begin{enumerate}[leftmargin=1.8em]
  \item 固定对手逐日计划，构造其“日期—起点—行动”时间表；
  \item 在状态 $(t,l,W,F)$ 上枚举移动、停留、挖矿，并按同边/同矿规则计算自身增量；
  \item 对相同状态只保留挖矿收入最高标签，用最短距离剪去无法到达状态；
  \item 交替求两人的最佳响应直至计划不变；
  \item 固定最终对手计划再完整求解一次，报告最大偏离增益。
\end{enumerate}
\end{resultbox}

\subsection{第一、二关完整逐日流水}

表\ref{tab:daily-ledger}由\path{output/problem1/known_weather_exact.json}自动生成；购买栏为“水/食物”。
\begin{longtable}{ccclrrrr}
\caption{第一、二关全局最优计划的逐日审计流水}\label{tab:daily-ledger}\\
\toprule
关卡 & 日 & 天气 & 位置/行动 & 现金 & 水 & 食物 & 购买(水/食)\\
\midrule
\endfirsthead
\multicolumn{8}{c}{\small 表\thetable\ （续）}\\
\toprule
关卡 & 日 & 天气 & 位置/行动 & 现金 & 水 & 食物 & 购买(水/食)\\
\midrule
\endhead
\midrule\multicolumn{8}{r}{续下页}\\\endfoot
\bottomrule\endlastfoot
1 & 0 & -- & 1/初始 & 5780 & 178 & 333 & 0/0\\
1 & 1 & 高温 & 25/行走 & 5780 & 162 & 321 & 0/0\\
1 & 2 & 高温 & 26/行走 & 5780 & 146 & 309 & 0/0\\
1 & 3 & 晴朗 & 23/行走 & 5780 & 136 & 295 & 0/0\\
1 & 4 & 沙暴 & 23/停留 & 5780 & 126 & 285 & 0/0\\
1 & 5 & 晴朗 & 21/行走 & 5780 & 116 & 271 & 0/0\\
1 & 6 & 高温 & 9/行走 & 5780 & 100 & 259 & 0/0\\
1 & 7 & 沙暴 & 9/停留 & 5780 & 90 & 249 & 0/0\\
1 & 8 & 晴朗 & 15/行走 & 4150 & 243 & 235 & 163/0\\
1 & 9 & 高温 & 13/行走 & 4150 & 227 & 223 & 0/0\\
1 & 10 & 高温 & 12/行走 & 4150 & 211 & 211 & 0/0\\
1 & 11 & 沙暴 & 12/挖矿 & 5150 & 181 & 181 & 0/0\\
1 & 12 & 高温 & 12/挖矿 & 6150 & 157 & 163 & 0/0\\
1 & 13 & 晴朗 & 12/挖矿 & 7150 & 142 & 142 & 0/0\\
1 & 14 & 高温 & 12/挖矿 & 8150 & 118 & 124 & 0/0\\
1 & 15 & 高温 & 12/挖矿 & 9150 & 94 & 106 & 0/0\\
1 & 16 & 高温 & 12/挖矿 & 10150 & 70 & 88 & 0/0\\
1 & 17 & 沙暴 & 12/停留 & 10150 & 60 & 78 & 0/0\\
1 & 18 & 沙暴 & 12/停留 & 10150 & 50 & 68 & 0/0\\
1 & 19 & 高温 & 12/挖矿 & 11150 & 26 & 50 & 0/0\\
1 & 20 & 高温 & 13/行走 & 11150 & 10 & 38 & 0/0\\
1 & 21 & 晴朗 & 15/行走 & 10470 & 36 & 40 & 36/16\\
1 & 22 & 晴朗 & 9/行走 & 10470 & 26 & 26 & 0/0\\
1 & 23 & 高温 & 21/行走 & 10470 & 10 & 14 & 0/0\\
1 & 24 & 晴朗 & 27/行走 & 10470 & 0 & 0 & 0/0\\
2 & 0 & -- & 1/初始 & 5300 & 130 & 405 & 0/0\\
2 & 1 & 高温 & 9/行走 & 5300 & 114 & 393 & 0/0\\
2 & 2 & 高温 & 18/行走 & 5300 & 98 & 381 & 0/0\\
2 & 3 & 晴朗 & 26/行走 & 5300 & 88 & 367 & 0/0\\
2 & 4 & 沙暴 & 26/停留 & 5300 & 78 & 357 & 0/0\\
2 & 5 & 晴朗 & 35/行走 & 5300 & 68 & 343 & 0/0\\
2 & 6 & 高温 & 36/行走 & 5300 & 52 & 331 & 0/0\\
2 & 7 & 沙暴 & 36/停留 & 5300 & 42 & 321 & 0/0\\
2 & 8 & 晴朗 & 37/行走 & 5300 & 32 & 307 & 0/0\\
2 & 9 & 高温 & 38/行走 & 5300 & 16 & 295 & 0/0\\
2 & 10 & 高温 & 39/行走 & 5200 & 10 & 283 & 10/0\\
2 & 11 & 沙暴 & 39/停留 & 3460 & 174 & 273 & 174/0\\
2 & 12 & 高温 & 47/行走 & 3460 & 158 & 261 & 0/0\\
2 & 13 & 晴朗 & 55/行走 & 3460 & 148 & 247 & 0/0\\
2 & 14 & 高温 & 55/挖矿 & 4460 & 124 & 229 & 0/0\\
2 & 15 & 高温 & 55/挖矿 & 5460 & 100 & 211 & 0/0\\
2 & 16 & 高温 & 55/挖矿 & 6460 & 76 & 193 & 0/0\\
2 & 17 & 沙暴 & 55/挖矿 & 7460 & 46 & 163 & 0/0\\
2 & 18 & 沙暴 & 55/挖矿 & 8460 & 16 & 133 & 0/0\\
2 & 19 & 高温 & 62/行走 & 4730 & 201 & 207 & 201/86\\
2 & 20 & 高温 & 55/行走 & 4730 & 185 & 195 & 0/0\\
2 & 21 & 晴朗 & 55/挖矿 & 5730 & 170 & 174 & 0/0\\
2 & 22 & 晴朗 & 55/挖矿 & 6730 & 155 & 153 & 0/0\\
2 & 23 & 高温 & 55/挖矿 & 7730 & 131 & 135 & 0/0\\
2 & 24 & 晴朗 & 55/挖矿 & 8730 & 116 & 114 & 0/0\\
2 & 25 & 沙暴 & 55/挖矿 & 9730 & 86 & 84 & 0/0\\
2 & 26 & 高温 & 55/挖矿 & 10730 & 62 & 66 & 0/0\\
2 & 27 & 晴朗 & 55/挖矿 & 11730 & 47 & 45 & 0/0\\
2 & 28 & 晴朗 & 55/挖矿 & 12730 & 32 & 24 & 0/0\\
2 & 29 & 高温 & 56/行走 & 12730 & 16 & 12 & 0/0\\
2 & 30 & 高温 & 64/行走 & 12730 & 0 & 0 & 0/0\\
\end{longtable}


\subsection{复现实验索引}

\begin{table}[H]
\centering
\caption{关键命令、输入与输出}
\begin{tabularx}{0.96\textwidth}{L{3.2cm} Y Y}
\toprule
命令 & 主要输入 & 主要输出\\
\midrule
\texttt{make test} & 场景、模拟器、策略与博弈代码 & 20项规则与回归测试\\
\texttt{make solve-known} & 第一、二关地图和30天天气 & \texttt{known\_weather\_exact.json}\\
\texttt{make analyze-problem2} & 第三、四关场景 & MDP敏感性和手工策略基线\\
\texttt{make optimize-level4} & 567候选、训练/测试种子 & 第四关搜索和20000条样本外结果\\
\texttt{make analyze-problem3} & 第五、六关多人场景 & 均衡、装载搜索、四策略比较\\
\texttt{make figures} & 全部JSON/CSV & 9张PDF/PNG候选图\\
\texttt{make result} & 精确流水、原Result模板 & 最终\texttt{Result.xlsx}\\
\texttt{make deliver} & 全部计算结果与论文源文件 & 最终PDF与Excel交付物\\
\bottomrule
\end{tabularx}
\end{table}

随机实验的训练、测试种子和样本数均写入对应JSON；所有输出表使用UTF-8 with BOM以便Excel直接打开。最终论文引用的图只从候选目录晋级到\path{paper/figures/generated/}，避免试验图混入交付。
